\section*{Chapitre \ref{ch:g3:separating-mixtures} --- Séparer les mélanges}

\begin{solution}{exo:g3:separating-mixtures:1}
Le \omterm{def:g3:separating-mixtures:sieve}{tamisage}~: les grains de riz sont plus petits que les petits pois et
passent à travers un tamis aux trous de la bonne taille. (Trier à la
main marche aussi, lentement.)
\end{solution}

\begin{solution}{exo:g3:separating-mixtures:2}
La terre descend lentement et forme une couche au fond~; l'eau au-dessus
s'éclaircit. C'est la \omterm{def:g3:separating-mixtures:decant}{décantation}.
\end{solution}

\begin{solution}{exo:g3:separating-mixtures:3}
Le \omterm{def:g3:separating-mixtures:filter}{filtrat} est le liquide qui traverse le filtre. Quand on fait du café,
le \omterm{def:g3:separating-mixtures:filter}{filtrat} est le café dans la verseuse~; le marc reste dans le papier.
\end{solution}

\begin{solution}{exo:g3:separating-mixtures:4}
Oui. Le sel est dissous, et un filtre ne retient pas ce qui est
dissous~: il passe avec l'eau.
\end{solution}

\begin{solution}{exo:g3:separating-mixtures:5}
En faisant évaporer l'eau~: au soleil, ou en chauffant doucement, l'eau
part dans l'air et le sucre reste dans la coupelle.
\end{solution}

\begin{solution}{exo:g3:separating-mixtures:6}
Le gravier reste dans le tamis. Les grains de sable sont plus petits que
les trous, alors ils passent à travers.
\end{solution}

\begin{solution}{exo:g3:separating-mixtures:7}
(a) Le \omterm{def:g3:separating-mixtures:sieve}{tamisage}. (b) La \omterm{def:g3:separating-mixtures:filter}{filtration} (ou la \omterm{def:g3:separating-mixtures:decant}{décantation} puis le
\omterm{def:g3:separating-mixtures:decant}{transvasement}). (c) L'\omterm{def:g3:separating-mixtures:evaporate}{évaporation}.
\end{solution}

\begin{solution}{exo:g3:separating-mixtures:8}
Quatre étapes~: les grilles, les bassins de \omterm{def:g3:separating-mixtures:decant}{décantation}, les filtres à
sable, les microbes tués. Ce sont les bassins de \omterm{def:g3:separating-mixtures:decant}{décantation} qui
enlèvent la boue.
\end{solution}

\begin{solution}{exo:g3:separating-mixtures:9}
Filtrer~; récupérer le sable sur le papier~; faire évaporer le \omterm{def:g3:separating-mixtures:filter}{filtrat}
pour obtenir le sel.
\end{solution}

\begin{solution}{exo:g3:separating-mixtures:10}
Non. Le sel de l'eau de mer est dissous, et un filtre ne retient pas ce
qui est dissous~: le \omterm{def:g3:separating-mixtures:filter}{filtrat} est toujours de l'eau salée.
\end{solution}

\begin{solution}{exo:g3:separating-mixtures:11}
Verser doucement le thé dans une tasse une fois les feuilles déposées
au fond (\omterm{def:g3:separating-mixtures:decant}{décantation} et \omterm{def:g3:separating-mixtures:decant}{transvasement}), ou le verser à travers une
passoire ou un filtre (\omterm{def:g3:separating-mixtures:sieve}{tamisage} ou \omterm{def:g3:separating-mixtures:filter}{filtration}).
\end{solution}
